\begin{abstract}
We study a quantum register whose Hilbert space grows: fresh qudits are born in a pure state, uncorrelated with the rest, and the register is scrambled by local dynamics. The observable is the deficit $G(A)=\log q\,\min(|A|,|\Abar|)-S(A)$ between the maximal and the actual entropy of a region. A velocity criterion in the spirit of Frautschi (relaxation versus expansion) predicts that the deficit closes when scrambling is fast enough; simulations refute it. We show that the controlling variable is causal, not kinetic.
(i)~For every unitary circuit, on any graph and with arbitrary births, $S(A)\le\log q\,|\Abar\cap\Inf(A)|$, where $\Inf(A)$ collects the qudits whose world line meets the causal past of $A$ (Theorem~A); a Lieb--Robinson version holds for local Hamiltonians (Theorem~B).
(ii)~For an isotropic front $R\propto t^\alpha$, the deficit is extensive for every local dynamics if $\alpha\ge1$, and can be made sub-extensive if $\alpha<1$ (Theorem~C); in one dimension the causal bound is attained.
(iii)~For gates drawn from a 2-design, the averaged purity of a prefix reduces \emph{exactly} to the walk of a single domain wall absorbed by the line of births. This yields $\E[G]=\Theta(m^{2-1/\alpha})$ for $\frac12<\alpha<1$ (Theorem~D) and, for one birth every $k$ layers, the exact annealed Rényi-2 density $\theta_k$ with $\frac{2q}{q^2+1}\cosh\theta_k=e^{-\theta_k/k}$ (Theorem~E).
(iv)~Stabilizer simulations up to $6400$ qubits and Haar state-vector simulations up to $22$ qubits agree with these statements; predictions were written before the computations, and the ones that failed are reported.
The deficit is robust to the choice of birth state and to local correlations with the past. In a free lattice field with a vacuum, however, uncorrelated births cost an energy of the order of the cutoff, so the cosmological reading of these results remains a constrained analogy.
\end{abstract}
